\section{Лекция 9}
\begin{lemma}
    Пусть $\Phi$ --- выпуклое подмножество в $\R^k$. В некоторой окрестности $U(\Phi)$ определена функция $\phi: U(\Phi)\to \R,\ \phi\in \mathcal{C}^1(\Phi)$. Тогда $\forall \overline{x}, \overline{y}\ \exists\ \overline{\xi}\in [\overline{x}, \overline{y}]$:
    \[\phi(\overline{x})-\phi(\overline{y})=\langle \grad \phi(\overline{\xi}, \overline{x}-\overline{y}) \rangle\]
\end{lemma}
\begin{proof}
    Пусть для $t\in [0,1]$: $h(t):=\phi(\overline{y}+t\cdot (\overline{x}-\overline{y}))$. Тогда 
    \begin{multline*}
        h^{\prime}(t)=(\phi(y_1+t(x_1-y_1),\dots,y_k+t(x_k-y_k)))'=\\
        =\sum\limits_{j=1}^{k}\phi_j'(\dots)\cdot (x_j-y_j)=\langle \grad \phi(\overline{y}+t(\overline{x}-\overline{y})), \overline{x}-\overline{y} \rangle
    \end{multline*}
\end{proof}
\noindentВспомним что такое норма в $\R^k$.
$\overline{x}=(x_1,\dots,x_k),\ \|\sqrt{x_1^2+\dots x_k^2}\|_2$
\begin{exercise} Доказать, что при $p\geq 1$
    $\|\overline{x}\|_p=(|x_1|^p+\dots+|x_k|^p)^{\frac{1}{p}}$
    является нормой.
\end{exercise}
\noindent$\|\overline{x}\|_1=(|x_1|^p+\dots+|x_k|^p)^{\frac{1}{p}}$ называется Манхэттенской нормой.\\
Также можно ввести $\|\overline{x}\|_{\infty}=\max\limits_{j} |x_j|$
\begin{lemma}
    Пусть $\Phi$ --- выпуклый компакт в $\R^k,\ U(\Phi)$ --- открытая окрестность $\Phi$. Пусть $\phi: U(\Phi)\to \R^m,\ \phi\in \mathcal{C}^1(\Phi)$. Тогда $\exists\ M>0\ \forall \overline{x}, \overline{y}\in \Phi$:
    \[\|\phi(\overline{x})-\phi(\overline{y})\|_{\R^m}\leq M\cdot \|\overline{x}-\overline{y}\|_{\R^k}\]
\end{lemma}
\begin{proof}
\begin{multline*}
    \sqrt{(\phi_1(\overline{x})-\phi_1(\overline{y}))^2+\dots+(\phi_m(\overline{x})-\phi_m(\overline{y}))^2}\leq \sum\limits_{i=1}^{m}|\phi_i(\overline{x})-\phi_i(\overline{y})|=\\
    \overset{(1)}{=}\sum\limits_{i=1}^{n}|\langle \grad \phi_i(\overline{\xi_i}),\ \overline{x}-\overline{y}\rangle|\overset{(2)}{=} \sum\limits_{i=1}^{m}\|\grad \phi_i(\overline{\xi}_i)\|_{\R_k}\cdot \|\overline{x}-\overline{y}|\|_{\R_k}
\end{multline*}
\[g_i(\overline(\xi)_i)=\sqrt{\sum\limits_{j=1}^{k}\left(\frac{\partial \phi_i}{\partial x_j}(\overline{\xi}_i)\right)^2}\]
$\exists\ M_i,\ \forall \overline{x}\in \Phi: g_u(\overline{x})\leq M_i$.
\end{proof}
Напоминание
Пусть $\Omega\subset \R^k,\ \phi: \Omega\to \R^m$. Матрицей Якоби называется мледующая матрица
\[\begin{pmatrix}
    \frac{\partial \phi_1}{\partial x_1} & \frac{\partial \phi_1}{\partial x_2} & \dots & \frac{\partial \phi_1}{\partial x_k}\\
    \frac{\partial \phi_2}{\partial x_1} & \frac{\partial \phi_2}{\partial x_2} & \dots & \frac{\partial \phi_2}{\partial x_k}\\
    \vdots & \vdots & \null & \vdots\\
    \frac{\partial \phi_m}{\partial x_1} & \frac{\partial \phi_m}{\partial x_2} & \dots & \frac{\partial \phi_m}{\partial x_k}\\
\end{pmatrix}\]
Если $m=k$, то определитель этой матрицы называют Якобианом, его обозначают $J$.
Пусть $\phi$ --- дифференцируема в точке $\overline{x}_0$. Тогда
\[\phi(\overline{x}_0+\Delta \overline{x})-\phi(\overline{x}_0)=J(\overline{x}_0)\cdot \Delta \overline{x}+\alpha(\Delta \overline{x})\cdot \|\Delta \overline{x}\|_{\R^k}\]
где $\alpha(\Delta \overline{x})=\bar{\bar{o}}(1)$.
\begin{lemma}
    Пусть $\Phi$ --- выпуклый компакт в $\R^k,\ U(\Phi)$ --- открытая окрестность $\Phi$. Пусть $\phi: U(\Phi)\to \R^m,\ \phi\in \mathcal{C}^1(\Phi)$. Тогда $\exists\ \omega(t): \R_+\to \R_+,\ \lim\limits_{t\to 0+}\omega(t)=0$ такая, что $\forall \overline{x},\ \overline{x}+\overline{\Delta x}\in \Phi$:
    \[\|\phi(\overline{x}+\overline{\Delta x})-\phi(\overline{x})-J(\overline{x})\cdot \overline{\Delta x}\|_{\R^m}\leq \omega(\|\overline{\Delta x}\|_{\R_k})\cdot \|\overline{\Delta x}\|_{\R^k}\] 
\end{lemma}
\begin{proof}
    \begin{multline*}
        \phi(\overline{x}+\Delta \overline{x})-\phi(\overline{x})=(\phi_1(\overline{x}+\Delta \overline{x})-\phi_1(\overline{x}), \dots, \phi_m(\overline{x}+\Delta \overline{x})-\phi_m(\overline{x}))=\\
        =(\langle \grad \phi_1(\overline{x}_1), \Delta \overline{x} \rangle, \dots, \grad \phi_m(\overline{x}_m), \Delta \overline{x})-(\langle \grad \phi_1(\overline{\xi_1}), \Delta \overline{x} \rangle, \dots, \grad \phi_m(\overline{\xi_m}), \Delta \overline{x})=\\
        =(\langle \grad \phi_1(\overline{x})-\grad \phi_1(\overline{\xi}_1, \Delta \overline{x}) \rangle, \dots, \langle \grad \phi_m(\overline{x})-\grad \phi_m(\overline{\xi}_m, \Delta \overline{x}) \rangle)
    \end{multline*}
    Тогда 
    \begin{multline*}
        \|\phi(\overline{x}+\Delta \overline{x})-\phi(\overline{x})\|\leq \sum\limits_{i=1}^{m}|\langle \grad \phi_i(\overline{x})-\grad \phi_i(\overline{\xi}_i), \Delta\overline{x} \rangle|\leq\\
        \leq \left(\sum\limits_{i=1}^{m}\|\grad \phi_i(\overline{x})-\grad \phi_i(\overline{\xi}_i)\|\right)\cdot \|\Delta \overline{x}_i\|
    \end{multline*}
    Введем 
    \[\omega_{ij}(\delta):=\sup\{\left|\frac{\partial \phi_i}{\partial x_j}(\overline{x})-\frac{\partial \phi_i}{\partial x_j}(\overline{y})\right|: \overline{x}, \overline{y}\in \Phi,\ \|\overline{x}-\overline{y}\|<\delta\}\]
    Теперь определим
    \[\omega(\delta)=\sum\limits_{i=1}^{m}\sum\limits_{j=1}^{k}\omega_{ij}(\delta)\]
    Ограничим первый множитель
    \[(1)\leq \sum\limits_{i=1}^{m}\sum\limits_{j=1}^{k}\left|\frac{\partial \phi_i}{\partial x_j}(\overline{x})-\frac{\partial \phi_i}{\partial x_j}(\overline{\xi}_i)\right|\leq \omega(\|\Delta \overline{x}\|)\]
\end{proof}
\begin{theorem}
    Пусть $\Omega, G$ --- области в $\R^k,\ \phi^\Omega\to G$ --- биекция, причем
    \begin{enumerate}
        \item $\phi\in \mathcal{C}^1(\Omega)$
        \item $\det J\neq 0$ на $\Omega$.
        \item $\Phi$ --- измеримый компакт.
    \end{enumerate}
    Тогда $\phi(\Phi)=K$ --- измеримый компакт в $G$, причем
    \[\mu(K)=\idotsint\limits_{\Phi}|\det{J}|\ dx_1\dots dx_k\]
\end{theorem}
\begin{example}
    ПРИМЕР (С КАРТИНКАМИ АХАХАХХА)
\end{example}